\documentclass[11pt]{article}
\usepackage[UTF8]{ctex}
\usepackage{amsmath,amssymb,amsthm}
\usepackage{etoolbox}

% 1. 自定义定理样式实现标题后强制换行
\newtheoremstyle{break}% 样式名称
  {\topsep}{\topsep}% 环境上下间距
  {\normalfont}% 正文字体
  {}% 缩进量
  {\bfseries}{.}% 头部字体和标点
  {\newline}% 头部后的间距：设置为 \newline 即可实现完美的换行效果
  {}% 头部时间/附加参数规范

% 应用新样式并声明常用的定理类环境
\theoremstyle{break}
\newtheorem{theorem}{定理}[section]
\newtheorem{lemma}[theorem]{引理}
\newtheorem{corollary}[theorem]{推论}
\newtheorem{definition}[theorem]{定义}

% 2. 利用 etoolbox 宏包动态修补 proof 环境，使其标题后自动换行
\patchcmd{\proof}{\item[\hskip\labelsep\itshape#1\@addpunct{.}]}{\item[\hskip\labelsep\itshape#1\@addpunct{.}]\hfill\par}{}{}

\begin{document}
固定重数之和集大小的多项式压缩

草稿

2026年4月29日

摘要

对固定的 $h$，令 $N(h,k)$ 为最小的 $N$，使得任一基数为 $k$ 的整数集所能达到的每一个基数 $|hA|$，也都能由 $\{1,\dots,N\}$ 的一个 $k$ 元子集达到。利用 Rajagopal 关于 $|hA|$ 可能取值范围的固定 $h$ 定理，我们证明了
\[N(h,k) \le k^{C_h}\]
对每个固定的 $h$ 成立。新的要素是用一个多项式直径的替代物来替换 Rajagopal 大值构造中的稀疏几何块。该替换是在低次 Hilbert 函数的层面上进行的，当 $h$ 固定时，这些函数是与 $h$ 重和集相关的唯一数据。
\section{引言}

对于一个有限集合 \(A \subset \mathbb{Z}\) 和一个正整数 \(h\)，记
\[hA = \{a_1 + \cdots + a_h : a_i \in A\}\]
允许重复。定义
\[R(h,k) = \{|hA| : A \subset \mathbb{Z}, |A| = k\}.\]
令 \(N(h,k)\) 是使得
\[R(h,k) = \{|hA| : A \subseteq \{1,\dots,N\}, |A| = k\}\]
成立的最小正整数 \(N\)。等价地，我们也可以在区间 \([0, N-1]\) 中进行讨论，并在最后进行平移。

本文的目的在于证明下述多项式压缩界。
\textbf{定理 1.1.} 对每个固定的 $h \geq 1$，存在一个常数 $C$，使得对所有 $k \geq 2$，有
\[ N(h,k) \leq k^{C h} h. \]
证明利用了关于可能的固定重数和集大小范围的 Rajagopal 定理。令
\[ M_{h,k} = hk - h + 1, \quad K_{h,k} = \binom{h+k-1}{h}, \]
并定义
\[ \Delta_{h,k} = \bigcup_{\ell=0}^{\min(h,k)-3} [M_{h,k} + \ell h + 1, M_{h,k} + \ell h + (h-2-\ell)]. \]
这里及下文中，$[u,v]$ 表示满足 $u \leq n \leq v$ 的整数 $n$ 的集合；若 $u > v$，则该区间为空集。
\textbf{定理 1.2} (Rajagopal [1]). 对于每个固定的 $h$，存在 $k_h$ 使得对于所有 $k > k_h$，有 \[ R(h,k) = [M_{h,k}, K_{h,k}] \setminus \Delta_{h,k}. \] 对于大的 $k$，该证明是构造性的，并且为 Rajagopal 范围内的所有值给出了多项式直径的证据。有限多个较小的 $k$ 值最后通过增加多项式界限中的常数而被吸收。

该范围的下半部分已由 Rajagopal 的填充集构造提供，其直径是多项式的。范围的上半部分是需要进行改变的地方。Rajagopal 的大值构造使用了稀疏几何级数；它们的低次生成函数具有正确的形状，但它们的直径在相关参数上是指数级的。我们用具有相同截断生成函数的多项式直径格点配件来替换它们。由于 $h$ 是固定的，一个固定维度的混合基映射然后将所得的格点集嵌入到 $\mathbb{Z}$ 中，且直径仅有多项式增长。
\section{生成函数与不交并}
在本笔记的其余部分，$h$ 是固定的。$O_h(\cdot)$、$\Theta_h(\cdot)$ 和 $\asymp_h$ 中隐含的常数可能依赖于 $h$，但不依赖于 $k$。

对于无挠交换群中的一个有限集 $A$，定义截断和集生成函数
\[
F_A(z) = \sum_{q=0}^{h} |qA| z^q.
\]
所有涉及 $F_A$ 的恒等式均在模 $z^{h+1}$ 意义下理解。
\textbf{定义 2.1}（无交并）\quad 设 $A \subset \mathbb{Z}^r$ 和 $B \subset \mathbb{Z}^s$ 是有限集。定义
\[ A \sqcup B = (A,0,1,0) \cup (0,B,0,1) \subset \mathbb{Z}^{r+s+2}. \]
迭代的无交并可以任意结合。
\textbf{引理 2.2.} 对于有限的 $A$ 和 $B$，有
\[F_{A \sqcup B}(z) \equiv F_A(z) F_B(z) \pmod{z^{h+1}}.\]
\textbf{证明} $q(A \sqcup B)$ 中的一个元素，其来自 $(A, 0, 1, 0)$ 的和项有唯一确定的个数 $i$，而来自 $(0, B, 0, 1)$ 的和项有 $q-i$ 个，因为最后两个坐标是 $(i, q-i)$。因此
\[
|q(A \sqcup B)| = \sum_{i=0}^{q} |iA| \cdot |(q-i)B|,
\]
这正是系数恒等式。
\textbf{引理 2.3 (混合进制嵌入).} 设 $A \subset \mathbb{Z}^D$ 的所有坐标都在 $[0, M]$ 内。若 $Q > hM$ 且
\[
\phi(x_1, \dots, x_D) = x_1 + Q x_2 + \cdots + Q^{D-1} x_D,
\]
那么 \(\phi\) 是 \(A\) 上的一个 \(h\) 阶 Freiman 同构。特别地，对每个 \(0 \leq q \leq h\)，有 \(|q\phi(A)| = |qA|\)，且
\[
\phi(A) \subset [0, D M Q D^{-1}].
\]
因此，若 \(D = O_h(1)\) 且 \(M \leq k^{O_h(1)}\)，则 \(\phi(A)\) 具有关于 \(k\) 的多项式直径。
\textbf{证明.} 若 A 中至多 h 个元素的和的每个坐标均落在区间 $[0, hM]$ 内。由于 $Q > hM$，\phi 下的像相等意味着其 $Q$-进制展开的各位数字相等，从而这些向量的和相等。
\section{多项式直径无和进位 gadget}
我们首先记录一个明晰的多项式直径的 B-集供应。
\textbf{引理 3.1}（多项式自由块）. 对于每个 $r \ge 1$，存在一个集合 $U \subset \mathbb{Z}^{h+1}$，其大小为 $r$，且坐标以 $rh$ 为界，使得对于 $0 \le q \le h$，有
\[
|qU| = \binom{r+q-1}{r}
\]
等价地，
\[
F_{U_r}(z) \equiv (1-z)^{-r} \pmod{z^{h+1}}.
\]
对于 $r=0$，我们取 $U=\emptyset$ 和 $F=1$.
\textbf{证明。} 取 $U = \{(1,i,i^2,\dots,i^h) : 1 \le i \le r\}$。 假设两个 $q$ 重和相等，且 $q \le h$。第一个坐标表明这两个和使用相同数目的 $q$ 项。其余坐标给出了两个指标多重集的第一个 $h$ 幂和的相等性。牛顿恒等式随后给出了初等对称函数的相等性，因此这两个多重集相等。从而所有 $q$ 重和都互不相同。 接下来，我们为 Rajagopal 证明中的两个稀疏几何块构造多项式替换。我们使用一个无关联基。
\textbf{定义 3.2.} 设 $u_1, \dots, u_r$ 是无挠交换群中的一个序列。若对任意满足 $c_i, c_i' \in \mathbb{N}_0$ 且 $\sum_i c_i, \sum_i c_i' \leq L$ 的系数，当
\[
\sum_i c_i u_i = \sum_i c_i' u_i
\]
时，能推出对所有 $i$ 有 $c_i = c_i'$，则称该序列是 \textbf{$L$-脱节的}。

在我们的应用中，取 $L = h^2$ 就足够了。具体地说，
\[
u_i = (1, i, i^2, \dots, i^{h^2}) \in \mathbb{Z}^{h^2+1} \quad (1 \leq i \leq r).
\]
与上面相同的牛顿恒等式论证表明这些向量是 $h^2$-脱节的，并且它们的坐标以 $r^{h^2}$ 为界。
\section{G-小工具}

对于 \(3 \leq m \leq h\) 和 \(r \geq 0\)，定义
\[
G_{m,r} = \{0\} \cup \{u_i, m u_i : 1 \leq i \leq r\}.
\]
因此 \(|G_{m,r}| = 2r+1\).
\textbf{引理 3.3.} 对 $3 \le m \le h$ 及 $r \ge 0$，有
\[
(1 - z^m)^r F(z) \equiv G_{m,r} (1 - z)^{2r+1} \pmod{z^{h+1}}.
\]
此外，定义
\[
P_{m,r}^G(z) = (1 - z^m)^r,
\]
则有
\[
P_{m,r}^G(z) - P_{m,r+1}^G(z) = z^m (1 - z^m)^r.
\]
\textbf{证明}. 考虑一个 $q$ 重和，其中 $q \leq h$。它具有形式
\[ 
\sum_i a_i u_i + \sum_i b_i (m u_i) = \sum_i (a_i + m b_i) u_i,
\]
其中 $a_i, b_i \geq 0$。由于 $\sum_i (a_i + m b_i) \leq m h \leq h^2$，系数向量 $c_i = a_i + m b_i$ 由和式决定。

将 $c_i$ 写为 $c_i = m B_i + A_i$，其中 $0 \leq A_i < m$。要得到此系数向量所需的非零加项的最小数目为 $\sum_i (A_i + B_i)$，而 $G_{m,r}$ 中的 $0$ 的副本可以将其填充至任意更大的长度。因此生成函数为
\[
\frac{1}{(1-z^m)^r} \left( \sum_{A=0}^{m-1} z^A \right)^r \left( \sum_{B \geq 0} z^B \right)^r = \frac{1}{(1-z)^{2r+1}}.
\]
最后的恒等式是直接的。
\section{H-组件}

对于 $2 \leq m \leq h-1$ 且 $r \geq 0$，定义
\[
H_{m,r} = \{u_i, m u_i : 1 \leq i \leq r\}.
\]
因此 $|H_{m,r}| = 2r$，且 $H_{m,0} = \emptyset$。
\textbf{引理 3.4.} 对于 $2 \le m \le h-1$ 及 $r \ge 0$，有
\[
\frac{P_{H_{m,r}}(z)}{F(z)} \equiv H_{m,r}(1-z)^{2r} \pmod{z^{h+1}},
\]
其中
\[
P_{H_{m,r}}(z) = \frac{(1-z^m)^r - z^{m-1} (1-z)^r}{1 - z^{m-1}}.
\]
此外，
\[
P_{H_{m,r}}(z) \equiv 1 - \binom{r}{2} z^{m+1} + \sum_{\nu=m+2}^{h} O_{h}(r^{\nu-m+1}) z^{\nu} \pmod{z^{h+1}},
\]
并且
\[
P_{H_{m,r}}(z) - P_{H_{m,r+1}}(z) \equiv r z^{m+1} + \sum_{\nu=m+2}^{h} O_{h}(r^{\nu-m}) z^{\nu} \pmod{z^{h+1}}.
\]
\textbf{证明}：一个来自 $H_{m,r}$ 的 $q$ 重和，满足 $q \leq h$，具有形式
\[ \sum_i a_i u_i + \sum_i b_i (m u_i) = \sum_i (a_i + m b_i) u_i. \]
同样地，系数向量 $c = a + m b$ 由和式决定，因为 $\sum_i c_i \leq m h \leq h^2$。
记
\[ c_i = m d_i + e_i, \quad 0 \leq e_i < m, \]
并令 $D = \sum_i d_i, E = \sum_i e_i$。对于一个固定的系数向量 $c$，选取满足 $0 \leq b_i \leq d_i$ 的 $b$，其长度为
\[ \sum_i (a_i + b_i) = m D + E - (m - 1) \sum_i b_i. \]
当 $b$ 取遍所有值时，和式 $\sum_i b_i$ 遍历从 $0$ 到 $D$ 的每一个整数。因此，这个系数向量恰好对以下次数有贡献：
\[ D + E, D + E + (m - 1), \dots, m D + E. \]
从而，
\[
F_{H_{m,r}}(z) \equiv \sum_{\substack{d_i \ge 0,\\ 0 \le e_i < m}} z^{D+E} \left(1 + z^{m-1} + \cdots + z^{(m-1)D}\right)
= \left(\frac{1 - z^m}{1 - z}\right)^r \frac{(1 - z)^{-r} - z^{m-1}(1 - z^m)^{-r}}{1 - z^{m-1}}
= \frac{(1 - z^m)^r - z^{m-1}(1 - z)^r}{(1 - z)^{2r}(1 - z^{m-1})}.
\]
这证明了 $P_{H_{m,r}}$ 的公式。
所述的展开式可通过将分子按 $z$ 的幂展开得到。对于一步差分，使用
\[
P_{H_{m,r+1}}(z) - P_{H_{m,r}}(z) = \frac{z^m\left( (1 - z^m)^r - (1 - z)^r \right)}{1 - z^{m-1}}.
\]
第一个非零项为 $r z^{m+1}$，且 $z^\nu$ 的系数为 $O(r \nu^{-m})$（对 $\nu \ge m+2$）。
\section{系数记号}
我们使用 Rajagopal 系数记号的截断版本。一个级数 $P(z) = \sum_{q=0}^h p_q z^q$ 属于 $\mathcal{O}_h[[Xz]]$，如果对于 $0 \leq q \leq h$ 有 $p_q = \mathcal{O}(X^q)$。它属于 $\Theta_{h,+}[[Xz]]$，如果对于所有 $0 \leq q \leq h$ 有 $p_q \asymp X^q$ 且 $p_q > 0$。
下面的基本控制引理是下文中唯一需要的分析估计。
\textbf{引理 4.1} 固定 $0 \le \alpha < 1$。设 $c \to \infty$，并设 $E_1(z), \dots, E_s(z)$ 为有界数量的截断级数，且满足
\[
E_i(z) = 1 + O[\![c^\alpha z]\!]_{h}.
\]
则
\[
\prod_{i=1}^{s} (1-z)^{-c} E_i(z) \in \Theta[\![c z]\!]_{h,+}.
\]
更一般地，若
\[
D(z) = a z^d (1 + O[[c \alpha z]])^h
\]
其中 \(a > 0\) 且 \(0 \leq d \leq h\)，则
\[
s \left( \prod_{i=1}^{h} D(z) (1-z)^{-c} E_i(z) \right) \in a z^d \Theta [[c z]].
\]
\textbf{证明}：只需证明第二个陈述；第一个陈述是 $a = 1$, $d = 0$ 的情形。对于 $q \ge d$，$z^q$ 的系数为
\[
\binom{c + q - d - 1}{q - d} a_h + O \left( a_h c^\alpha \sum_{e=1}^{q-d} c^{q-d-e} \right).
\]
主项为 $\asymp a_h c^{q-d}$，而每个误差项为 $O \left( a_h c^{q-d-(1-\alpha)e} \right) = o \left( a_h c^{q-d} \right)$。因此，对于所有充分大的 $c$，该系数为正且 $\asymp a_h c^{q-d}$。
\section{范围的下部}

我们引用 Rajagopal 的填充集引理，以一种记录了直径的形式表述。
\textbf{定义 5.1.} 设 $B \subseteq [0,d]$。若满足 $hB = [0,hd]$，则称 $B$ 是 \textbf{$(h,d)$-填充}的。
\textbf{引理 5.2（Rajagopal 填充引理 [1]）}  
对于固定的 $h \geq 2$ 和所有充分大的 $k$，若
\[
\frac{(k-1)h}{k-2} \leq d \leq \frac{4(h^2+h)}{2},
\]
则存在一个 $(h,d)$-填充集 $B \subseteq [0,d]$，使得 $|B| = k-1$ 且 $[0, k/8] \subseteq B$。
\textbf{命题 5.3 (小值具有多项式见证)}. 对于固定的 $h \geq 2$，存在常数 $a > 0$ 和 $C > 0$，使得对于所有充分大的 $k$，每一个
$t \in [M_{h,k}, a k^h] \setminus \Delta_{h,k}$
都等于 $|hA|$，其中 $A \subseteq [0, C k^h]$ 且 $|A| = k$。
\textbf{证明}：令 \[ D_{\max} = \left\lceil \frac{(k-1)h}{4 \binom{h+2}{h} / 2} \right\rceil. \] 设 $d \in [k-2, D_{\max}]$，$i \in [1, h]$，且 $B \subseteq [0, d]$ 为由引理 5.2 提供的集合。对于充分大的 $k$，包含关系 $[0, k/8] \subseteq B$ 蕴含着 $[0, h^2] \subseteq B$。定义 \[ A = \{0\} \cup (i + B). \] 则 $|A| = k$ 且 $A \subseteq [0, h + D_{\max}]$。由于 $hB = [0, hd]$ 且 $[0, h^2] \subseteq B$，有 \[ hA = \{0\} \cup [i, h(i+d)]. \]
确实，区间 $[i, hi]$ 包含在 $i + B$ 中，而区间 $[hi, h(i+d)]$ 等于 $hi + hB$。因此，
\[
|hA| = h(i+d) - i + 2.
\]
令 $i+d = k + \ell$，则
\[
|hA| = hk + h\ell - i + 2 = M_{h,k} + h\ell + (h - i + 1).
\]
对于固定的 $\ell$，当 $i$ 取遍所有容许值时，这恰好给出了第 $\ell$ 行在 $\Delta_{h,k}$ 之外的整数，即
\[
M_{h,k} + h\ell + [\max(1, h-\ell-1), h].
\]
当 $d$ 的范围上至 $D_{\max}$ 时，这覆盖了 $\Delta_{h,k}$ 之外直到某个 $a_h kh$ 的所有值，其中 $a_h > 0$ 仅依赖于 $h$。其直径为 $O_h(kh)$。
\section{大部分范围}

我们现在证明 Rajagopal 大值构造的多项式直径替换。
\subsection{稠密块}
我们需要与 Rajagopal 相同的稠密块。对于 \(b \ge 3\)，令
\[
s = (h-1)(b-2)+1,
\]
且对于 \(1 \le j \le s\) 定义
\[
B_{j,b} = [0, b-2] \cup \{h(b-2)+2-j\}.
\]
因此 \(|B_{j,b}| = b\).
\textbf{引理 6.1} (Rajagopal 的稠密块 [1]). 对于固定的 $h$ 和 $b \ge 3$，集合 $B_{j,b}$ 满足：
\begin{enumerate}
\item[(a)] 对于 $2 \le \eta \le h$ 和 $2 \le j \le s_b$，
\[ 0 \le |\eta B_{j-1,b}| - |\eta B_{j,b}| \le h; \]
\item[(b)] 对于 $1 \le \eta \le h$ 和 $1 \le j \le s_b$，
\[ |\eta B_{j,b}| = O_h(b); \]
\item[(c)] 
\[ F_{B_{1,3}}(z) \equiv (1-z)^{-3} \pmod{z^{h+1}}; \]
\item[(d)] 对于 $b > 3$，
\[ F_{B_{1,b}}(z) \equiv (1-z)^{-1} F_{B_{s_{b-1},b-1}}(z) \pmod{z^{h+1}}. \]
\end{enumerate}
\section{稀疏多项式块}
固定指数
\[ \alpha = \frac{9}{10}, \quad \beta = \frac{4}{5}, \quad \gamma = \frac{7}{10}. \]
它们的选择仅需满足
\[ 0 < \beta < \alpha < 1, \quad 2\beta > 1, \quad 2\beta\gamma > 1. \]
固定 \( b \in [3, k - \lfloor k\gamma \rfloor] \) 并令 \( c = k - b \)。记
\[ S = \lfloor c^{\alpha} \rfloor, \quad T = \lfloor c^{\beta} \rfloor. \]
对参数
\[ 0 \le r_m \le S \ (3 \le m \le h), \quad 1 \le u_m \le T \ (2 \le m \le h-1), \]
定义
\[ R = c - \sum_{m=3}^{h} (2r_m + 1) - \sum_{m=2}^{h-1} 2u_m. \]
对于充分大的 \( k \)，这是非负的。定义稀疏块
\[ C(\mathbf{r}, \mathbf{u}) = \bigsqcup_{m=3}^{h} \mathcal{G}_{m, r_m} \sqcup \bigsqcup_{m=2}^{h-1} \mathcal{H}_{m, u_m} \sqcup \mathcal{U}_R. \]
那么 \( |C(\mathbf{r}, \mathbf{u})| = c \)，并且\textbf{引理 2.2}、\textbf{引理 3.1}、\textbf{引理 3.3} 和 \textbf{引理 3.4} 给出
\[ \frac{1}{(1-z)^c} F_{C}(z) \equiv \prod_{m=3}^{h} (1 - z^m)^{r_m} \prod_{m=2}^{h-1} P_{\mathcal{H}_{m, u_m}}(z) \pmod{z^{h+1}}. \tag{1} \]
\textbf{引理 6.2} (稀疏块的变体). 对固定的 $h$ 和充分大的 $c$，以下结论成立。

\textbf{(a)} 若 $3 \leq \mu \leq h$，则将 $r$ 增加 $1$ 将使 $F(z)$ 减小一个属于 $\mu C z^\mu \Theta[[cz]]_{h,+}$ 的元素。

\textbf{(b)} 若 $2 \leq \mu \leq h-1$，则将 $u$ 增加 $1$ 将使 $F(z)$ 减小一个属于 $\mu C u z^{\mu+1} \Theta[[cz]]_{\mu h,+}$ 的元素。
\textbf{证明.} 所有乘积均截断模 $z^{h+1}$。由 (1) 中余下的因子均具有形式
$1+O [[c\alpha z]]$，因为 $r \le c\alpha$ 且 $u \le c\beta \le c\alpha$。
\[ h m m \]
对于 (a)，由 \textbf{引理 3.3}，分子中的正下降为
\[ z^\mu(1-z^\mu)r_\mu = z^\mu(1+O [[c\alpha z]]). \]
乘以其他分子因子以及 $(1-z)^{-c}$，\textbf{引理 4.1} 便给出了所声称的
\[ h \]
$z^\mu\Theta [[cz]]$ 下降。
\[ h,+ \]
对于 (b)，\textbf{引理 3.4} 给出
\[ P_{H}(z)-P_{H}(z) = u z^{\mu+1}(1+O [[c\alpha z]]) \]
\[ \mu,u\mu \qquad \mu,u\mu+1 \qquad \mu \]
直到次数 $h$。再次应用 \textbf{引理 4.1}。
\section{连接性引理}

我们使用 Rajagopal 的区间引理，此处为完备性而将其列出。
\textbf{引理 6.3.} 设 $f : [0,n_1] \times \dots \times [0,n_d] \to \mathbb{Z}$ 在每个坐标上非增。对于 $1 \le i \le d$，定义
\[
\delta_i = \max \big( f(\mathbf{x}) - f(\mathbf{x} + \mathbf{e}_i) \big),
\]
其中最大值取遍所有满足 $x_i < n_i$ 的 $\mathbf{x}$，并定义
\[
\Delta_i = \min \big( f(x_1, \dots, x_{i-1}, 0, x_{i+1}, \dots, x_d) - f(x_1, \dots, x_{i-1}, n_i, x_{i+1}, \dots, x_d) \big).
\]
若 $\delta_1 \le 1$ 且 $\delta_i \le \Delta_{i-1}$ 对每一个 $2 \le i \le d$ 成立，则 $f$ 的像是一个整数区间。
\textbf{证明}：对于固定的 $x_2, \dots, x_d$，改变 $x_1$ 会得到一个区间，因为逐次下降量是 $0$ 或 $1$。

归纳假设，在改变 $x_1, \dots, x_{i-1}$ 且固定后续坐标之后，其像是一个区间。当 $x_i$ 增加 $1$ 时，新区间的起点至多比旧区间低 $\delta_i$，而旧区间的长度至少为 $\Delta_{i-1}$。因此，连续的区间会互相重叠。对 $i$ 进行归纳即可证明该命题。
\subsection{固定 b：连通性}
对于固定的 $b$ 和如上参数，定义
\[ A_{j,b} = B_b \sqcup C(r,u), \qquad 1 \leq j \leq s. \]
设 $\mathbb{A}_b$ 为由所有这样的 $A$ 构成的族。
\textbf{引理 6.4.} 对于固定的 $h \ge 3$ 和所有充分大的 $k$, 集合
\[\{|hA| : A \in \mathcal{A}\}\]
对于每个 $b \in [3, k - \lfloor k \gamma \rfloor]$ 都是一个整数区间。
\textbf{证明} 将变量排序为
\[ r_h, u_{h-1}, r_{h-1}, u_{h-2}, \dots, r_3, u_2, j. \]
共有 $2h-3$ 个变量。我们对相应的盒子上的函数 $f = |hA|_{j,b}$ 应用 \textbf{引理 6.3}，平移变量使得每个范围从 0 开始。

因为 $A = B \sqcup C$，由 \textbf{引理 2.2} 可得
\begin{equation}
|hA|_{j,b} = \sum_{\eta=0}^{h} |\eta C| \, |(h-\eta) B|. \label{eq:2}
\end{equation}

由 \textbf{引理 6.2} 和 \textbf{引理 6.1(a)}，函数 $f$ 关于每个变量均非增。将 $r_h$ 增加 1 仅改变次数 $h$ 中的 $F$，其中它使 $|hC|$ 恰好减少 1；因此 $\delta_C = 1$。
接下来需要比较单步下降量 \(\delta\) 与全变差 \(\Delta\)。我们记录所需的估计。若变量为 \(r_\mu\) 且 \(3 \le \mu \le h\)，由引理 6.2 及 (2) 式给出
\[
\delta(r_\mu) \le C_h \left( c^{h-\mu} + b c^{h-\mu-1} \right),
\]
其中负指数项已省略。由于每一步都出现相同的正主导项，在全范围内改变 \(r\) 可得
\[
\Delta(r_\mu) \ge c S_h \left( c^{h-\mu} + b c^{h-\mu-1} \right).
\]
类似地，若变量为 \(u_\mu\) 且 \(2 \le \mu \le h-1\)，则
\[
\delta(u_\mu) \le C_h T \left( c^{h-\mu-1} + b c^{h-\mu-2} \right),
\]
且
\[
\Delta(u_\mu) \ge c T^2 \left( c^{h-\mu-1} + b c^{h-\mu-2} \right).
\]
事实上，总下降量是权重为 \(u_\mu = 1,2,\dots,T-1\) 的单步下降量之和，其和为 \(\asymp T^2\)。最后，由引理 6.1(a) 及 (2)，将 \(j\) 改变 1 仅改变稠密块因子，有
\[
\delta(j) \le C_h |\eta C| \le C_h c^{h-2} \sum_{\eta=0}^{h-2} 1 = C_h c^{h-2}.
\]
现在比较相邻变量。由于 \(\beta < \alpha\)，方程 (5) 和 (4) 表明，对充分大的 \(c\)，每个 \(u_\mu\)-下降量至多为前一个 \(r_{\mu+1}\) 全变差。对于反向比较，注意到 \(c \ge k\gamma + O(1)\) 且 \(b \le k\)，因此条件 \(2\beta\gamma > 1\) 给出 \(b \le c^{2\beta - o(1)}\)。结合 \(2\beta > 1\)，方程 (3) 和 (6) 表明每个 \(r_\mu\)-下降量至多为前一个 \(u_\mu\) 全变差。最后，由于 \(2\beta > 1\)，(7) 式被 \(\mu=2\) 时的 (6) 式控制。

因此引理 6.3 的假设成立，且 \(f\) 的像是一个区间。
\subsection{变化 \(b\)}
\textbf{命题 6.5}（大值存在多项式见证）。固定 $h \ge 3$ 和 $\varepsilon > 0$。对所有充分大的 $k$，每个满足 $t \in [\varepsilon k^{h}, K]_{h,k}$ 的整数都等于某个包含在长度为 $k^{O_h(1)}$ 的区间中的 $k$ 元集 $A$ 对应的 $|hA|$。
\textbf{证明}：令
\[
A' = \bigcup_{b=3}^{k-\lfloor k\gamma \rfloor} \mathcal{A}_b.
\]
根据\textbf{引理 6.4}，每个 $\mathcal{H}_{\mathcal{A}_b} := \{ |h_A| : A \in \mathcal{A}_b \}$ 是一个区间。

我们首先证明相邻区间是重叠的。对于固定的 $b > 3$，在 $\mathcal{A}_b$ 中取具有极小稀疏参数以及 $j = 1$ 的元素。其稀疏块具有生成函数 $(1-z)^{-c}$。
在 $\mathcal{A}_{b-1}$ 中取具有极小稀疏参数以及 $j = s_{b-1}$ 的元素。其稀疏块具有生成函数 $(1-z)^{-(c+1)}$。由\textbf{引理 6.1(d)} 可得
\[
F_{B_{1,b}}(z)(1-z)^{-c} \equiv F_{B_{s_{b-1},b-1}}(z)(1-z)^{-(c+1)} \pmod{z^{h+1}}.
\]
因此，所选的两个集合具有相同的 $|hA|$ 值，且 $hA_b \cap hA_{b-1} \neq \emptyset$。故 $hA$ 是一个区间。

在 $b = 3$ 处，取所有稀疏参数为极小值且 $j = 1$，\textbf{引理 6.1(c)} 给出
\[F_A(z) \equiv (1-z)^{-3}(1-z)^{-(k-3)} = (1-z)^{-k} \pmod{z^{h+1}}.\]
因此
\[|hA| = \binom{h+k-1}{h} = K_{h,k}.\]

在另一端，取 $b = k - \lfloor k^\gamma \rfloor$，将所有稀疏参数设为极大值，并取 $j = s_b$。则 $c = \lfloor k^\gamma \rfloor$。对于 $A$ 中的每个元素，我们有 $|\eta C| = O_h(c^\eta)$ 且对于 $q \ge 1$ 有 $|qB_{j,b}| = O_h(b)$，而 $|0B_{j,b}| = 1$。因此该端点有
\[|hA| \le C_h (b^{ch-1}) = O_h(k^{\gamma h + 1 + \gamma(h-1)}) = o_h(k^h).\]

因此，对所有充分大的 $k$，区间 $hA$ 包含 $[\varepsilon k^h, K_{h,k}]$。

最后只需记录半径。稠密块 $B_{j,b}$ 位于一个长度为 $O_h(k)$ 的区间内。自由块 $U$ 的坐标以 $kh$ 为界。用于 $G_{m,r}$ 和 $H_{m,u}$ 的基的坐标以 $k^{O_h(1)}$ 为界，因为 $r, u \le k$。仅有 $O_h(1)$ 个不交并分量，因此所得格集位于固定维数中，且坐标以 $k^{O_h(1)}$ 为界。\textbf{引理 2.3} 将其嵌入到 $\mathbb{Z}$ 中，半径为 $k^{O_h(1)}$，并保持所有 $q \le h$ 的 $q$ 重和集大小不变。
\section{2 重情况}
我们对 \(h=2\) 的情形给出初等的二次论证，从而使 \textbf{定理 1.1} 在未涵盖于大值构造的端点情形下是自洽的。
\textbf{命题 7.1.} 对于每个 $k \geq 1$，有
\[N(2,k) \leq 8k^2 + 2k + 2.\]
\textbf{证明}。只需在 $[0, L]$ 中讨论，并在最后平移 $1$ 即可。令
\[
K = \binom{n+1}{2}, \quad \Delta_n = K - (2n - 1) = \binom{n-1}{2}.
\]
每个 $k$ 元集的二重和集的规模介于 $2k-1$ 与 $K_k$ 之间。反过来，固定 $t \in [2k-1, K_k]$，并设 $D = K_k - t$。选取最小的 $m$ 使得 $D \leq \Delta_m$，并令
\[
s = \Delta_m - D.
\]
若 $m = 1$，取 $B = \{0\}$。若 $m \geq 2$，则 $0 \leq s \leq m-2$，并取
\[
B = \{0, 1, \dots, m-2\} \cup \{m-1 + s\}.
\]
记 $L = \max B$，则 $L \leq 2m$ 且
\[
|B + B| = K_m - D.
\]
对于 $m = 1$ 这是直接的。对于 $m \ge 2$，若 $P = \{0, \dots, m-2\}$ 且 $x = m-1+s$，则 $P+P = [0, 2m-4]$，$x+P = [x, x+m-2]$，这些区间重叠或接触，且 $2x$ 贡献一个额外的点，得到 $|B+B| = 2m-1+s = K_m - D$。

令 $r = k-m$。若 $r = 0$，取 $A = B$ 且计数已经正确。若 $r = 1$，取 $A = B \cup \{2L+1\}$。则三个部分 $B+B$、$B+\{2L+1\}$ 和 $\{4L+2\}$ 是不交的，所以 $|A+A| = |B+B| + m + 1 = K_{m+1} - D = t$。因此我们可以假设 $r \ge 2$。选择一个奇素数 $p$ 满足 $r \le p < 2r$，并设
\[ Q = \max(2p, L+p) + 1, \quad c_i = Qi + (i^2 \bmod p), \quad 0 \le i < r. \]
集合 $C = \{c_0, \dots, c_{r-1}\}$ 是 Sidon 集。事实上，$c_i + c_j = c_u + c_v$ 首先给出 $i+j = u+v$，因为 $Q > 2p-2$，然后模 $p$ 约化给出 $i^2 + j^2 \equiv u^2 + v^2$，从而 $ij \equiv uv$ 且 $\{i, j\} = \{u, v\}$。又 $c_{i+1} - c_i > L$，所以平移 $B+c_i$ 是不交的。令
\[ M = c_{r-1} + L + 1, \quad S = \{M + c_i : 0 \le i < r\}, \quad A = B \cup S. \]
则 $B+B$、$B+S$ 和 $S+S$ 是不交的。此外 $|B+S| = mr$ 且 $|S+S| = K_r$。因此
\[ |A+A| = (K_m - D) + mr + K_r = K_k - D = t. \]
最后，因为 $p < 2r$ 且 $L \le 2m$，
\[ \max A \le 2r(L+4r+1) + L + 1 \le 8k^2 + 2k + 1. \]
平移 $1$ 后，这个集合位于 $[1, 8k^2+2k+2]$ 中。
\begin{proof}
我们现在完成\textbf{定理 1.1} 的证明。情形 $h = 1$ 是平凡的。情形 $h = 2$ 即为\textbf{命题 7.1}。

固定 $h \geq 3$。设 $a$ 如\textbf{命题 5.3} 中所给，并利用 $\varepsilon = a/2$ 应用\textbf{命题 6.5}。对于所有充分大的 $k$，\textbf{命题 5.3} 在一个多项式长度的区间内实现了
\[[M, a k^h] \setminus \Delta_{h,k}\]
中的所有值，同时\textbf{命题 6.5} 在一个多项式长度的区间内实现了
\[[(a/2)k^h, K_{h,k}]\]
中的所有值。因此
\[[M, K] \setminus \Delta_{h,k}\]
中的所有值都有一个多项式直径的见证。由 Rajagopal 的值域定理（\textbf{定理 1.2}），这恰好就是对于所有充分大的 $k$ 的 $R(h,k)$。

对于有限多个较小的 $k$ 值，有限性是直接的：对于每个可达的值选择一个见证集，将其平移至正整数中，并取这有限多个选择中的最大端点。如果需要的话增加指数 $C$，这些有限多个情形也被限制在 $k^{C_h}$ 之内。最后，上述所有构造都是在非负区间内进行的。通过 $+1$ 平移将该集合置于 $\{1,\dots,N\}$ 中，并且不改变 $|hA|$。
\end{proof}
\textbf{注记 8.1 (下界)} 最大值 $K = \binom{h+k-1}{h}$ 由一个 $B_h$-集取得。如果 $A \subseteq [1,N]$ 满足 $|hA| = K_{h,k}$，那么 $hA \subseteq [h, hN]$，从而
\[
\binom{h+k-1}{h} \leq hN - h + 1.
\]
因此
\[
N(h,k) \gg k^h.
\]
从而多项式上界在 $k$ 的指数增长阶中具有正确的定性阶，直到指数的值。
\end{document}